% Préambule — Documentation utilisateur Tockem SPE
\usepackage[utf8]{inputenc}
\usepackage[T1]{fontenc}
\usepackage[french]{babel}
\usepackage{geometry}
\geometry{a4paper, margin=2.5cm}
\usepackage{graphicx}
\usepackage{booktabs}
\usepackage{array}
\usepackage{longtable}
\usepackage{hyperref}
\usepackage{fancyhdr}
\usepackage{float}
\usepackage{xcolor}
\usepackage{ifthen}

\hypersetup{
    colorlinks=true,
    linkcolor=blue!60!black,
    urlcolor=blue!60!black,
    pdftitle={Documentation utilisateur Tockem SPE},
    pdfauthor={Erick Tsafack}
}

\pagestyle{fancy}
\fancyhf{}
\fancyhead[L]{\small Volet EHA TOCKEM}
\fancyhead[R]{\small Documentation utilisateur Tockem SPE}
\fancyfoot[C]{\thepage}
\renewcommand{\headrulewidth}{0.4pt}

% Inclusion d'une capture d'écran (images/ relative au dossier latex/)
% Usage : \inclurecapture{nom-fichier.png}{Légende}{fig:label}
\newcommand{\inclurecapture}[3]{%
    \begin{figure}[H]
        \centering
        \IfFileExists{images/#1}{%
            \includegraphics[width=0.92\textwidth]{images/#1}%
        }{%
            \fcolorbox{gray!50}{gray!10}{%
                \parbox[c][3.5cm][c]{0.88\textwidth}{%
                    \centering\small
                    \textbf{Capture d'écran à insérer}\\[0.4em]
                    \texttt{images/#1}%
                }%
            }%
        }%
        \caption{#2}
        \label{#3}
    \end{figure}%
}
